\documentclass[UTF8,a4paper,11pt]{ctexart}
\usepackage[margin=2.2cm]{geometry}
\usepackage{tikz}
\usetikzlibrary{arrows.meta,positioning,fit,shapes.geometric}
\usepackage{graphicx}
\usepackage{booktabs}
\usepackage{array}
\usepackage{longtable}
\usepackage{xcolor}
\usepackage{hyperref}
\hypersetup{colorlinks=true,linkcolor=blue!45!black,urlcolor=blue!45!black}
\setlength{\emergencystretch}{3em}

\definecolor{iaBlue}{HTML}{EAF3FF}
\definecolor{iaGreen}{HTML}{ECF7EF}
\definecolor{iaRed}{HTML}{FFF0EE}
\definecolor{iaGold}{HTML}{FFF7DF}
\definecolor{iaGray}{HTML}{F3F5F8}
\definecolor{iaInk}{HTML}{172033}

\tikzset{
  box/.style={draw=iaInk, line width=.8pt, rounded corners=4pt, align=center, minimum height=1.1cm, text width=3.1cm, inner sep=5pt},
  bluebox/.style={box, fill=iaBlue, draw=blue!55!black},
  greenbox/.style={box, fill=iaGreen, draw=green!45!black},
  redbox/.style={box, fill=iaRed, draw=red!55!black},
  goldbox/.style={box, fill=iaGold, draw=orange!70!black},
  graybox/.style={box, fill=iaGray, draw=gray!65!black},
  arr/.style={-{Latex[length=2.5mm]}, line width=.75pt, draw=iaInk},
  darr/.style={arr, dashed}
}

\title{\heiti WealthCraftAgent 股票分析 Agent 架构设计文档}
\author{2352048 夏浩博 \quad 2351269 黄一和}
\date{\today}

\begin{document}
\maketitle

\section{项目概述}
WealthCraftAgent 是一个基于 .NET 10、WPF 与 Microsoft Semantic Kernel 的股票分析 Agent。项目面向 A 股分析场景，输入股票代码或股票名称后，系统从主营业务、K 线结构、财务报表、新闻动态四个角度生成分析，并支持缠论 RAG、历史相似走势 RAG 与会话内多轮追问。

本项目覆盖课程要求的五个核心要素：通过 OpenAI 兼容接口集成 LLM；通过 \texttt{InvestAgentLoop} 和会话编排实现 Thought--Action--Observation 推理循环；通过 Semantic Kernel 注册多个自定义工具；通过 \texttt{ConversationMemory} 与会话状态保存上下文；通过 WPF 桌面端提供用户交互界面。

\section{总体架构}
系统采用“桌面 UI + 核心 Agent 编排 + 工具/数据服务”的分层结构。WPF 负责输入、图表、历史记录与流式展示；Core 层负责 Agent 会话、子 Agent 调度、Semantic Kernel 调用、数据抓取和本地知识检索。

\begin{figure}[htbp]
\centering
\resizebox{\textwidth}{!}{
\begin{tikzpicture}[
  node distance=1.1cm and 1.25cm,
  box/.style={draw=iaInk, line width=.85pt, rounded corners=4pt, align=center, minimum height=1.0cm, text width=3.0cm, inner sep=5pt},
  bluebox/.style={box, fill=iaBlue, draw=blue!55!black},
  greenbox/.style={box, fill=iaGreen, draw=green!45!black},
  redbox/.style={box, fill=iaRed, draw=red!55!black},
  goldbox/.style={box, fill=iaGold, draw=orange!70!black},
  graybox/.style={box, fill=iaGray, draw=gray!65!black},
  arr/.style={-{Latex[length=2.5mm]}, line width=.75pt, draw=iaInk},
  darr/.style={arr, dashed}
]
  \node[bluebox] (ui) {WPF 桌面端\\输入、图表、多轮对话\\步骤可观测};
  \node[goldbox, right=of ui, text width=3.45cm] (a) {Agent A 编排层\\会话创建、任务规划\\状态合成、最终报告};
  \node[greenbox, right=of a, yshift=2.65cm] (b) {Agent B\\K 线结构\\缠论 RAG\\历史走势 RAG};
  \node[greenbox, right=of a] (c) {Agent C\\公司新闻\\行业新闻\\事件情绪归因};
  \node[greenbox, right=of a, yshift=-2.65cm] (d) {Agent D\\财务指标\\盈利 / 成长\\负债趋势};
  \node[graybox, right=1.35cm of c, text width=3.9cm, minimum height=3.35cm] (tools) {工具与数据层\\StockPrice / MarketNews\\FinancialReport / TechnicalAnalysis\\LocalKnowledge\\东方财富 / Yahoo / Finnhub\\缠论文档与历史案例库};
  \node[redbox, below=1.2cm of tools, text width=3.9cm] (reply) {回复反馈\\子 Agent 分析结果\\综合结论与风险提示};
  \node[bluebox, below=1.2cm of a, text width=3.45cm] (mem) {Memory 与持久化\\ConversationMemory\\SessionState\\SQLite 分析历史};

  \draw[arr] (ui) -- (a);
  \draw[arr] (a) -- (b);
  \draw[arr] (a) -- (c);
  \draw[arr] (a) -- (d);
  \draw[arr] (b.east) -- (tools.west);
  \draw[arr] (c.east) -- (tools.west);
  \draw[arr] (d.east) -- (tools.west);
  \draw[arr] (tools) -- (reply);
  \draw[arr] (reply.west) .. controls +(-2.2,-.7) and +(2.0,-.55) .. (mem.east);
  \draw[darr] (mem.north) -- (a.south);
\end{tikzpicture}
}
\caption{WealthCraftAgent 多 Agent 股票分析架构图}
\end{figure}

\section{核心组件设计}
\subsection{LLM 集成}
项目使用 Microsoft Semantic Kernel 统一管理 LLM 调用。依赖注入在 \texttt{AddInvestAgent} 中完成：读取 \texttt{AgentOptions} 中的模型、接口地址、API Key 和代理配置，创建 Kernel，并通过 \texttt{AddOpenAIChatCompletion} 接入 OpenAI 兼容接口。

LLM 调用分为两类：一类是 \texttt{InvestAgentLoop} 中的自动工具选择，用自动函数选择让模型决定调用哪个 Kernel Function；另一类是多 Agent 编排中的 \texttt{AgentPromptRunner}，负责流式输出、重试、推理标签过滤和图像输入降级。架构图中不把 LLM 单独画成业务节点，而是把它视为工具调用与回复生成过程中的基础能力：B/C/D 先连接到工具与数据层，工具层产生可用于回复的观察结果，再写回 Memory 与持久化模块供 Agent A 继续编排。

\subsection{Agent Loop 与会话编排}
项目同时保留了通用 ReAct 循环和更适合业务演示的确定性多 Agent 工作流。

\begin{itemize}
  \item \texttt{InvestAgentLoop} 实现课程要求的 Thought--Action--Observation 循环：每轮先产生 Thought 步骤，调用 LLM，若返回工具调用则执行 Action 并记录 Observation，若无工具调用则输出最终 Response。
  \item \texttt{SessionAnalysisOrchestrator} 作为 Agent A，负责股票会话创建、初始分析并行调度、追问意图识别、子 Agent 派单和最终报告合成。
  \item Agent B/C/D 分别处理 K 线、新闻、财务，职责单一，便于扩展和答辩解释。
\end{itemize}

\begin{figure}[htbp]
\centering
\resizebox{\textwidth}{!}{
\begin{tikzpicture}[
  node distance=1.0cm and 1.05cm,
  box/.style={draw=iaInk, line width=.85pt, rounded corners=4pt, align=center, minimum height=1.02cm, text width=2.85cm, inner sep=5pt},
  bluebox/.style={box, fill=iaBlue, draw=blue!55!black},
  greenbox/.style={box, fill=iaGreen, draw=green!45!black},
  redbox/.style={box, fill=iaRed, draw=red!55!black},
  goldbox/.style={box, fill=iaGold, draw=orange!70!black},
  graybox/.style={box, fill=iaGray, draw=gray!65!black},
  arr/.style={-{Latex[length=2.5mm]}, line width=.75pt, draw=iaInk},
  looparr/.style={-{Latex[length=2.5mm]}, line width=.8pt, draw=purple!70!black, dashed}
]
  \node[bluebox] (input) {用户输入\\股票代码 / 股票名称\\或会话追问};
  \node[graybox, right=of input] (session) {会话准备\\解析标的\\写入 Memory};
  \node[goldbox, right=of session] (thought) {Thought\\理解问题\\判断所需信息};
  \node[bluebox, right=of thought] (action) {Action\\调用工具\\或执行 B/C/D};
  \node[greenbox, below=1.45cm of action, minimum height=1.18cm] (obs) {Observation\\接收数据\\状态 Patch 回写};
  \node[goldbox, below=1.05cm of obs, text width=3.1cm] (condition) {条件判断\\信息是否足够回答};
  \node[redbox, left=1.45cm of condition] (response) {Response\\生成分析结论\\风险提示};
  \node[graybox, left=of response] (save) {历史保存\\会话记录\\支持回放};
  \node[goldbox, above=.9cm of action, text width=2.85cm] (guard) {保护机制\\MaxSteps\\异常处理};

  \draw[arr] (input) -- (session);
  \draw[arr] (session) -- (thought);
  \draw[arr] (thought) -- (action);
  \draw[arr] (action) -- (obs);
  \draw[arr] (obs) -- (condition);
  \draw[arr] (condition) -- node[above, font=\small] {信息充分} (response);
  \draw[arr] (response) -- (save);
  \draw[looparr] (condition.north west) -- node[pos=.55, left, font=\small] {信息不足} (thought.south);
  \draw[arr] (guard) -- (action);
\end{tikzpicture}
}
\caption{Agent 推理流程图}
\end{figure}

\section{工具与数据设计}
\begin{longtable}{p{3.2cm}p{4.2cm}p{7.2cm}}
\toprule
插件 / 服务 & 主要能力 & 设计说明\\
\midrule
\texttt{StockPrice} & 当前价格、历史日 K、月 K、股票搜索、资金流 & 为 Agent B 和 UI 图表提供行情基础数据；组合数据源优先适配 A 股。\\
\texttt{FinancialReport} & 关键财务指标、利润分析 & 为 Agent D 提供 ROE、ROA、毛利率、净利率、增长率、负债率等序列。\\
\texttt{MarketNews} & 新闻抓取、市场情绪 & 为 Agent C 提供公司新闻、行业新闻和情绪线索。\\
\texttt{Technical} \texttt{Analysis} & MA、RSI、MACD、交易信号 & 支持通用技术指标计算，也可被 ReAct 循环自动调用。\\
\texttt{LocalKnowledge} & 缠论文档检索、缠论图片检索 & 从 \texttt{docs/} 读取本地知识库，为缠论 RAG 提供文本片段和图例资源。\\
\bottomrule
\end{longtable}

数据服务层使用组合模式：\texttt{CompositeStockDataService} 汇聚东方财富、Yahoo、Alpha Vantage、Finnhub 等来源，并配合 \texttt{FileHttpCache} 降低重复请求。历史相似走势 RAG 由 \texttt{HistoricalPatternService} 基于本地案例库提取特征并匹配类似结构，输出样本数量、相似原因和后续收益分布，避免把历史案例误写成确定性预测。

\section{记忆、状态与可观测性}
\texttt{ConversationMemory} 保存系统提示、用户消息、助手回复和工具返回，超过最大轮次时自动保留最近上下文，避免提示词无限膨胀。会话内结构化状态由 \texttt{AnalysisSessionState} 维护，包括股票信息、K 线缓存、新闻列表、财务序列、各子 Agent 结果、最终报告等。

可观测性通过 \texttt{AgentStep} 实现。每个 Thought、Action、Observation、Response 都会通过 \texttt{IAnalysisStreamingObserver} 推送到 WPF，用户能够看到 Agent 正在抓取数据、检索知识库、生成分析或处理异常。

\section{异步、错误处理与扩展性}
项目大量使用 \texttt{async/await}：股票数据、新闻、财务指标、LLM 流式输出、初始 B/C/D 并行分析均为异步流程。初始分析中，Agent A 对 B/C/D 使用 \texttt{Task.WhenAll} 并发执行，提高首轮报告生成速度。

错误处理分为三层：LLM 调用失败时返回用户可读错误；工具或子 Agent 失败时记录日志并跳过该模块继续分析；数据不足时在提示词和最终报告中显式说明不确定性。扩展新工具时，只需实现新的服务或 Kernel 插件，并在 DI 中注册；扩展新子 Agent 时，实现 \texttt{ISubAgentService} 即可纳入编排。

\section{课程要求对应关系}
\begin{tabular}{p{3.6cm}p{10.2cm}}
\toprule
课程要求 & 项目实现\\
\midrule
LLM 集成 & Semantic Kernel + OpenAI ChatCompletion，支持流式输出和自动函数选择。\\
推理循环 & \texttt{InvestAgentLoop} 实现 ReAct；\texttt{SessionAnalysisOrchestrator} 实现业务化多 Agent 工作流。\\
工具调用 & 注册 5 个插件，多于课程要求的 3 个自定义工具。\\
Memory & \texttt{ConversationMemory}、\texttt{AnalysisSessionState} 与 SQLite 历史记录共同维护上下文。\\
用户界面 & WPF 桌面端，包含 K 线图、新闻、财务、聊天和历史回放。\\
加分项 & 多 Agent 协作、RAG、流式输出、可观测步骤、单元测试。\\
\bottomrule
\end{tabular}

\end{document}
